\section{Mina arbetsuppgifter under LIA-perioden}
\label{sec:arbetsuppgifter}

Under LIA 2 fick jag ansvar för flera infrastrukturförbättringar som alla hade direkt påverkan
på Insightas produktionsmiljö. Arbetet berörde fem tydliga områden: flytt av
orkestreringssystemet till Kubernetes, säkerhet och IAM, driftstabilitet vid deployments,
strukturering av container-arkitekturen och uppdatering av CI/CD-pipelinen.

GitHub-arbetet visar också tydligt hur mycket som faktiskt gjordes under perioden. I detta repo
skapade jag totalt 21 pull requests. De var inte slumpmässiga ändringar, utan en serie mindre
steg som tillsammans byggde upp en större lösning. Först lade jag grunden med kluster,
behörigheter, Artifact Registry och RBAC. Därefter arbetade jag med Helm, Dockerfile,
loggning och GitHub Actions för deployment. När den grunden fanns på plats gick arbetet vidare
till stabilitet, till exempel att separera tjänster i Docker Compose, låta jobb fortsätta under
deployment och minska onödiga omstarter. Den sista delen handlade om underhåll, som att hålla
CI/CD-flödet uppdaterat när GitHub Actions gick över till Node.js 24.

Det här arbetssättet blev viktigt för mig. I stället för att försöka leverera en stor lösning i en
enda pull request arbetade jag i små steg. Det gjorde arbetet lättare att granska, lättare att
testa och lättare att ändra när något behövde förbättras.

\subsection{Flytta datapipelinesystemet till Kubernetes i Google Cloud}
\label{subsec:k8s}

Insightas system för att schemalägga och köra datapipelines körde tidigare på en virtuell
maskin med Docker. Det fungerade, men lösningen var svår att skala och deployment krävde
flera manuella steg. Min uppgift var att flytta systemet till ett Kubernetes-kluster i Google
Cloud (GKE) och göra deployment mer stabil och mer automatiserad.

Jag byggde \textbf{Helm-mallar} för att paketera och deploya systemet och dess databas på ett
enhetligt sätt. Konfigurationen delades upp i ett baslager med gemensamma inställningar och
ett lager för respektive miljö. Det gör det enklare att hantera både staging och produktion
utan att behöva kopiera samma konfiguration flera gånger.

Jag använde \textbf{Terraform CDKTF} för att skapa och hantera resurser i Google Cloud som
kod. Det gällde bland annat IAM-roller, service accounts, Artifact Registry för
container-images och klustrets maskintyper. På så sätt blev infrastrukturen versionshanterad i
Git och lättare att återskapa.

Jag satte också upp ett \textbf{automatiserat deploymentflöde i GitHub Actions} som deployar
systemet till klustret vid releaser. Flödet skapar ett eget Kubernetes-namespace för varje
release-branch och rensar det automatiskt när branchen stängs. Det gör staging-miljön mer
ordnad och minskar manuellt arbete.

\subsection{Säker identitet -- Workload Identity istället för statiska nycklar}
\label{subsec:security}

En viktig del av projektet var hur Kubernetes-pods autentiserar sig mot Google Cloud-tjänster.
Det gamla sättet byggde på att en JSON-nyckel från ett GCP service account sparades i ett
Kubernetes secret. Det är en säkerhetsrisk eftersom en läckt nyckel kan användas av någon som
inte ska ha åtkomst.

Jag införde \textbf{Workload Identity}, vilket är ett säkrare sätt att lösa autentisering i GKE.
Med den lösningen kopplas en Kubernetes Service Account direkt till ett GCP Service Account.
Podden får då rätt identitet utan att någon nyckel behöver sparas i klustret. För mig var detta
ett tydligt exempel på hur teorin om minsta möjliga behörighet blev praktisk verklighet.

Samtidigt konfigurerade jag IAM-roller med minsta möjliga behörighet och satte upp RBAC i
Kubernetes via Terraform. Det gjorde att även säkerhetsinställningarna blev tydliga,
versionshanterade och lätta att följa upp.

\subsection{Driftstabilitet -- deployment utan avbrott i pågående jobb}
\label{subsec:stability}

Ett konkret driftsproblem var att varje gång ny kod deployades stoppades pågående
datapipeline-körningar. Det berodde på att hela Docker Compose-stacken togs ned och
startades om. För ett system som kör schemalagda jobb är det ett allvarligt problem, eftersom
ett misslyckat jobb kan försena viktig data för hela organisationen.

Lösningen blev att låta varje jobb köra i sin \textbf{egen container}, fristående från
kärnservicarna. Då kan man starta om webbgränssnittet och andra permanenta delar utan att
störa ett jobb som redan körs. Jag testade detta i praktiken genom att starta ett längre jobb,
stänga ned övriga containers och sedan starta upp dem igen. Jobbet slutfördes utan fel.

\subsection{Strukturering av container-arkitekturen}
\label{subsec:compose}

En bakomliggande orsak till stabilitetsproblemet var att flera olika delar av systemet delade
samma container. Det gjorde lösningen svår att underhålla och svår att felsöka.

Jag refaktorerade därför Docker Compose-konfigurationen så att varje komponent körs i sin
\textbf{egen container}. Det gav tre tydliga fördelar: det blev enklare att skala enskilda delar,
enklare att läsa loggar och enklare att göra mer kontrollerade deployments.

Jag förbättrade också deploymentskriptet så att det bara återskapar de containers som verkligen
behöver det. Kärnservicar startas om bara om de är trasiga eller stoppade. Kod-containern
återskapas alltid vid ny deployment så att ny kod verkligen används. Gamla överblivna
containers rensas bort automatiskt. Det gjorde deployment snabbare och mer förutsägbar.

\subsection{Uppdatering av CI/CD-pipelinen}
\label{subsec:cicd}

GitHub Actions, som vi använder för all CI/CD-automatisering, slutade stödja äldre versioner
av Node.js i sina actions. Därför behövde flera workflows uppdateras så att pipelinen inte
skulle sluta fungera.

Jag gick igenom workflows, hittade de actions som behövde uppdateras och bytte till versioner
som fungerar med Node.js 24. Jag tog också bort \texttt{latest}-taggar. Det är en liten men
viktig förbättring, eftersom en obestämd version kan ändras utan varning och skapa fel i
pipelinen.

\subsection*{Sammantagen bild}

Alla uppgifter under LIA 2 hänger ihop kring ett gemensamt tema: att göra ett produktionssystem
\textit{säkrare}, \textit{mer stabilt} och \textit{mer automatiserat}. Arbetet täcker samtliga
kärnområden i LIA 2-kursmålen -- infrastruktur som kod, CI/CD, containerorkestrering,
säkerhet och automatisering -- och har haft direkt påverkan på teamets dagliga arbete.

\subsection*{Koppling till utbildningen}

Arbetet under LIA 2 kopplar tydligt till flera delar av utbildningen. Inom molntjänster använde
jag GCP och GKE i praktiken. Inom IT-säkerhet arbetade jag med IAM, RBAC och säkrare
autentisering. Inom operativsystem och nätverk märktes kunskaperna i Docker, Linux-miljöer,
deploymentflöden och felsökning. Det gjorde att sådant vi läst i skolan fick ett tydligt
praktiskt sammanhang.